\section{The Cost of Post-Quantum Migration}
\label{sec:pq}

The analysis in this section requires no computation and no measurement. Every
figure follows arithmetically from the wire-format layout of
Section~\ref{sec:modelling} and the published parameter sizes of
ML-KEM~\cite{fips203}. That makes it the most robust part of the paper: it
does not depend on the machine it was produced on, and it cannot be disturbed
by a prover's resource limits.

\subsection{Packet growth}

\begin{table}[t]
\centering
\caption{In-band KEM substitution, sizes in bytes. The substitution modelled
is the cheapest one possible, so these are a lower bound on any in-band
design.}
\label{tab:pqsizes}
\footnotesize
\begin{tabular}{@{}lrrrr@{}}
\toprule
KEM & $ek$ & $ct$ & initiation & response \\
\midrule
X25519 (baseline) & 32 & 32 & 148 & 92 \\
\midrule
ML-KEM-512 & 800 & 768 & 916 ($\times$6.2) & 828 ($\times$9.0) \\
ML-KEM-768 & 1184 & 1088 & 1300 ($\times$8.8) & 1148 ($\times$12.5) \\
ML-KEM-1024 & 1568 & 1568 & 1684 ($\times$11.4) & 1628 ($\times$17.7) \\
\bottomrule
\end{tabular}

\end{table}

We model the cheapest possible in-band substitution: the initiator's 32-byte
ephemeral public value becomes an encapsulation key, the responder's becomes
the corresponding ciphertext, and every other field keeps its length. Because
this is the minimal construction, the resulting sizes
(Table~\ref{tab:pqsizes}) are a \emph{lower bound} on any in-band design,
which is what makes the conclusion below robust rather than contingent on a
particular protocol sketch. An initiation grows by a factor of \wgKemGrowth{}
at ML-KEM-768 and \wgKemTenGrowth{} at ML-KEM-1024.

\subsection{Path reachability}

\begin{figure}[t]
\centering
\resizebox{\columnwidth}{!}{\begin{tikzpicture}[x=1cm,y=1cm]
  \footnotesize
  \fill[okgreen!78] (0,-0.22) rectangle (0.608,0.22);
  \node[anchor=east] at (-0.12,-0.0) {X25519 (baseline)};
  \node[anchor=west] at (0.728,-0.0) {\bfseries 148\,B};
  \fill[okgreen!78] (0,-0.97) rectangle (3.766,-0.53);
  \node[anchor=east] at (-0.12,-0.75) {ML-KEM-512};
  \node[anchor=west] at (3.886,-0.75) {\bfseries 916\,B};
  \fill[badred!78] (0,-1.72) rectangle (5.344,-1.28);
  \node[anchor=east] at (-0.12,-1.5) {ML-KEM-768};
  \node[anchor=west] at (5.464,-1.5) {\bfseries 1300\,B};
  \fill[badred!78] (0,-2.47) rectangle (6.923,-2.03);
  \node[anchor=east] at (-0.12,-2.25) {ML-KEM-1024};
  \node[anchor=west] at (7.043,-2.25) {\bfseries 1684\,B};
  \draw[badred,dashed,thick] (5.065,-2.75) -- (5.065,0.78);
  \node[badred,anchor=east,inner sep=2pt,font=\scriptsize,font=\scriptsize\bfseries] at (5.065,0.78) {IPv6 minimum MTU: 1232\,B};
  \draw[black!55,dashed,thin] (6.052,-2.75) -- (6.052,1.22);
  \node[black!55,anchor=west,inner sep=2pt,font=\scriptsize] at (6.052,1.22) {Ethernet, IPv4: 1472\,B};
  \draw[->] (0,-3.1) -- (7.400,-3.1) node[anchor=west] {\scriptsize bytes};
  \draw (0.000,-3.2) -- (0.000,-3.0) node[anchor=north,yshift=-4pt] {\scriptsize 0};
  \draw (2.467,-3.2) -- (2.467,-3.0) node[anchor=north,yshift=-4pt] {\scriptsize 600};
  \draw (4.933,-3.2) -- (4.933,-3.0) node[anchor=north,yshift=-4pt] {\scriptsize 1200};
  \draw (7.400,-3.2) -- (7.400,-3.0) node[anchor=north,yshift=-4pt] {\scriptsize 1800};
\end{tikzpicture}
}
\caption{Initiation size against the usable UDP payload of each path. The
ML-KEM-768 bar crosses the IPv6 minimum-MTU rule; that crossing is the whole
of the argument in this section.}
\label{fig:mtu}
\end{figure}

\begin{table}[t]
\centering
\caption{ML-KEM-768 initiation against \wgPaths{} representative path MTUs.
Only the first row is a guarantee; the others are common cases.}
\label{tab:pqmtu}
\footnotesize
\begin{tabular}{@{}lrrr@{}}
\toprule
Path & MTU & usable UDP & ML-KEM-768 initiation \\
\midrule
IPv6 minimum MTU & 1280 & 1232 & \textbf{exceeds by 68\,B} \\
Ethernet, IPv4 & 1500 & 1472 & fits, 172\,B spare \\
Ethernet, IPv6 & 1500 & 1452 & fits, 152\,B spare \\
PPPoE, IPv4 & 1492 & 1464 & fits, 164\,B spare \\
Tunnel stacking, IPv4 & 1400 & 1372 & fits, 72\,B spare \\
\bottomrule
\end{tabular}

\end{table}

RFC~8200~\cite{rfc8200} requires every IPv6 link to carry a
\wgIpvsixMtu-byte datagram. After an IPv6 header of 40\,B and a UDP header of
8\,B, the usable payload is \wgIpvsixPayload\,B. An in-band ML-KEM-768
initiation is \wgKemInit\,B. It exceeds that payload by \wgOvershoot\,B, and a
datagram that does not fit the IPv6 minimum MTU has no guaranteed path
(Figure~\ref{fig:mtu}, Table~\ref{tab:pqmtu}). Path MTU discovery for datagram
transports~\cite{rfc8899} can find a larger path when one exists, but it
cannot manufacture one where it does not.

\para{The chain, stated plainly.} Each link below is individually
unremarkable, and the conclusion is not:

\begin{quote}
KEM parameter sizes \(\rightarrow\) datagram size \(\rightarrow\) path MTU
\(\rightarrow\) fragmentation or PMTUD failure \(\rightarrow\) the handshake
is unreachable on some paths \(\rightarrow\) the architecture must change.
\end{quote}

This is why post-quantum designs for WireGuard either restructure the
handshake around KEMs from the start~\cite{hulsing2021pqwireguard} or perform
the key agreement over a separate channel and inject the result into the
pre-shared key slot. Rosenpass~\cite{rosenpass2023} takes the second route,
and the analysis above supplies the reason it is necessary rather than merely
convenient.

\para{What migration costs beyond bytes.} There is a real tension worth
naming. WireGuard's single-packet, stateless handshake is itself an
engineering asset: it needs no fragment reassembly, holds no state before
authentication, and is what makes the cookie mechanism possible. Post-quantum
migration threatens precisely that asset. ML-KEM-512 would fit every path
considered at \wgKemFiveInit\,B, but at a security level most deployments will
not choose; ML-KEM-1024, at \wgKemTenInit\,B, fits none of them.

\subsection{Three dimensions, one conclusion}

Combining this section with Section~\ref{sec:results}:

\begin{itemize}\itemsep2pt
\item \emph{Packet size.} An in-band initiation grows by a factor of
      \wgKemGrowth{} at ML-KEM-768 and \wgKemTenGrowth{} at ML-KEM-1024.
\item \emph{Path reachability.} Above ML-KEM-512, in-band substitution loses
      the guarantee of a usable path.
\item \emph{Assumption strength.} A KEM-derived pre-shared key preserves the
      properties of a configured one only if the KEM binds its shared secret
      to the encapsulation key; without that binding the pre-shared key is
      recoverable while every primitive remains unbroken (P12a/P12b).
\end{itemize}

The cost of post-quantum migration is not merely computational. It appears
along all three dimensions at once, and a design that addresses only one of
them has not addressed the problem.
